\chapter{Server HTTP e applicazioni server-side}

\section{Il Server HTTP}

\begin{definizione}[Server HTTP]
Un \textbf{server HTTP} \`e un'applicazione eseguita su una macchina detta \emph{host}, che
sfrutta (fra le altre cose) i servizi TCP/IP del sistema operativo per ricevere richieste
HTTP e rispondervi. Dall'esterno ci si rivolge all'host con l'indirizzo IP (o lo hostname
corrispondente); per contattare lo specifico servizio serve la \textbf{porta} su cui \`e in
esecuzione. Il suo compito, in breve: restare in ascolto sulla porta assegnata e, all'arrivo
di una richiesta, elaborarla e fornire una risposta.
\end{definizione}

La maggior parte dei server HTTP \`e in esecuzione sulla porta di default \textbf{80}
(\textbf{443} per HTTPS), che quindi non va esplicitata. In fase di \emph{sviluppo e test} le
cose cambiano: i servizi comunicano sullo stesso computer -- IP \texttt{127.0.0.1}, hostname
\texttt{localhost} (lo strato TCP/IP sa che destinatario = mittente) -- e l'unico modo per
differenziarli \`e la porta: ecco perch\'e spesso si hanno pi\`u server HTTP su porte diverse.

\section{Gestione delle richieste: due modelli}

Un server HTTP deve gestire una grande quantit\`a di richieste ``in parallelo'': le richieste
giungono in \emph{coda} e vengono considerate in ordine di arrivo, ma il tempo di
elaborazione T \`e potenzialmente superiore all'intervallo fra due arrivi (la coda \`e
raramente vuota). Non si pu\`o gestire la coda in modo \emph{completamente} sequenziale --
sarebbe come un laboratorio di analisi che non accetta un paziente finch\'e tutti quelli che
lo precedono non hanno completato gli esami e sono usciti: spreco di risorse per chi eroga e
di tempo per chi riceve.

\begin{importante}[Modello 1: single-thread asynchronous event loop]
Le richieste sono prese in carico in modo sequenziale da un \textbf{loop a thread singolo};
qualsiasi elaborazione necessaria deve quindi avvenire \textbf{in modo asincrono}:
\begin{itemize}
  \item le elaborazioni di \emph{I/O} (file system, altri servizi come un DBMS) sono
        naturalmente asincrone: le esegue il sistema operativo, non il processo del server;
  \item le elaborazioni \emph{CPU-intensive} vanno gestite generando un thread separato.
\end{itemize}
Analogia: al banco dell'accettazione c'\`e \emph{una sola persona}, che esamina la richiesta
e la redirige: se \`e un'attesa (stampa di un referto) avvia l'operazione e chiede di
attendere; se richiede attivit\`a lunghe chiama un collega che prende in carico l'utente. In
entrambi i casi passa subito alla richiesta successiva.
Esempi: \textbf{Nginx}, qualsiasi server basato su \textbf{Node.js}.
\end{importante}

\begin{importante}[Modello 2: multi-thread]
A ogni richiesta il server genera un \textbf{thread separato e dedicato}, che la prende in
carico e pu\`o gestirla anche in modo sincrono. Analogia: ogni volta che un utente entra,
appare un nuovo sportello con una persona interamente dedicata a lui.
Esempi: \textbf{Apache Web}, \textbf{Apache Tomcat}.
\end{importante}

A un esame superficiale il multi-thread sembrerebbe pi\`u costoso ma pi\`u efficiente verso
gli utenti; in realt\`a negli ultimi due decenni sono emersi server single-thread persino
pi\`u efficienti (\textbf{Nginx \`e praticamente nato per dimostrarlo}). Un vantaggio del
single-thread: \textbf{performance pi\`u prevedibili sotto alto carico} -- spesso \`e meglio
una performance leggermente peggiore ma prevedibile (si pensi all'attesa alla fermata
dell'autobus!). Con le ottimizzazioni attuali i due approcci sono comparabili; resta
importante \emph{sapere su quale modello si basa} il server su cui si ``monta'' la propria
applicazione.

\section{HTTP Server vs Web Server}

\begin{importante}
Il \textbf{Web} \`e un servizio specifico realizzato tramite HTTP: un server \`e un
\emph{Web Server} se realizza davvero il ``core business'' -- fornire risorse statiche e
dinamiche che contribuiscono al WWW, in modo sufficientemente general purpose (di solito si
dice che un applicativo ``\emph{fa da}'' Web Server, perch\'e dipende da come \`e usato; in
linea di principio potrebbe esistere un Web Server non basato su HTTP\dots\ servirebbe
per\`o il browser corrispondente).
\textbf{HTTP} \`e invece un elemento \emph{infrastrutturale}: ``Server HTTP'' si riferisce al
protocollo con cui comunica, non agli scopi. Un server HTTP pu\`o non essere usato come Web
Server (es. esporre i dati di un DB, fare solo da reverse proxy); un'applicazione server-side
che espone un'API HTTP non \`e certo un Web Server, ma \emph{contiene} necessariamente un
server HTTP per ricevere le richieste.
\end{importante}

\section{Le funzioni di un Server Web/HTTP}

\subsection{Fornire contenuti statici e dinamici}
La funzione pi\`u comune (non l'unica): fornire contenuti statici e dinamici via HTTP --
inclusa la possibilit\`a per gli utenti di \emph{inviare} dati da pubblicare (la base del
``Web 2.0'') e le vere Applicazioni Web. Le risorse scambiate sono principalmente JSON, ma
un'app client-side ha sempre come punto d'ingresso una pagina HTML + CSS + JS: il server che
la ``pubblica'' deve fornire queste risorse (statiche o dinamiche); quasi sempre un minimo di
gestione di risorse statiche serve.

\subsection{Virtual hosting}
Un Web Server pu\`o ospitare \textbf{pi\`u host}, ognuno col proprio hostname, avendo (nella
maggior parte dei casi) un \emph{unico indirizzo IP}: tutti gli hostname virtuali sono
mappati dal DNS sullo stesso IP, quindi il server deve \textbf{esaminare la richiesta HTTP}
per stabilire quale hostname era stato usato. Alternativa: mappare ogni hostname su una
\emph{porta} diversa -- ma il DNS non mappa la porta e il browser usa quella di default:
serve allora il \textbf{reverse proxying}.
\begin{definizione}[Reverse proxy]
Un \textbf{reverse proxy} (per HTTP) \`e un server HTTP che riceve le richieste (es. sulla
porta di default) e \emph{in base al loro contenuto le reindirizza altrove} (es. al server
principale sulla porta abbinata allo hostname della richiesta). Molto usato anche per
siti/app i cui servizi sono erogati da \emph{molteplici} Web Server: stesso hostname, e il
proxy smista in base al pathname. Alcuni server (come Nginx) sono usati principalmente per
questo.
\end{definizione}
(Nota: in una sottorete locale si pu\`o usare un IP come alias di un altro, dando a ogni host
virtuale il proprio IP; \`e per\`o meno efficiente, e si preferisce mappare pi\`u hostname
sullo stesso IP.)

\subsection{Autorizzazione e autenticazione}
Molti Web Server permettono o vietano l'accesso a certi pathname/risorse in base
all'\emph{IP del mittente} (es. rete aziendale: risorse interne solo ai computer della rete,
sito pubblico per tutti; i dipendenti in smart working usano una \textbf{VPN}, che permette
di ``presentarsi'' con un IP della rete locale -- legalmente, installando un certificato del
proprietario della rete).
\begin{attenzione}[Non sono la stessa cosa]
\textbf{Autenticazione} = \emph{riconoscere l'identit\`a} dell'utente che richiede una
risorsa. \textbf{Autorizzazione} = \emph{permettere o vietare l'accesso} a una risorsa: pu\`o
dipendere da molti fattori, e l'identit\`a \`e solo uno di essi. Spesso usate in sinergia
(per proteggere contenuti privati), ma distinte.
\end{attenzione}

\subsection{Caching}
Il server pu\`o memorizzare le risposte HTTP associate alle relative richieste, per
\emph{evitare di reperire di nuovo la stessa risorsa} a fronte della stessa richiesta --
tanto pi\`u se la risorsa \`e dinamica e costosa da calcolare, o va richiesta a un altro
server (tipicamente un reverse proxy fa anche caching). Obiettivi: velocizzare la risposta e
ridurre il carico. La risposta non pu\`o per\`o essere ``la stessa in eterno'': uso,
validit\`a e scadenza della cache si \textbf{negoziano fra client e server tramite gli header
HTTP}.

\subsection{Bandwidth throttling}
Quando il server eroga dati che consumano molta banda in uscita, pu\`o \emph{ridurre
artificialmente} l'ampiezza di banda di una singola risposta, per suddividere la banda fra
pi\`u stream paralleli (analogia dei rubinetti su un tubo comune). Es.: video streaming in
alta qualit\`a -- senza throttling, il primo stream occupa tutta la banda e il secondo utente
va ``a singhiozzo''.

\subsection{Rewrite engine}
Non sempre la URL richiesta dal client corrisponde al pathname interno con cui il server
reperisce la risorsa -- spesso \`e voluto: URL interne lunghe e complesse, ma verso i client
URL ``pulite'' ed evocative. Alcuni server includono un \textbf{motore di riscrittura}: regole
sintattiche che traducono le URL dalla versione ``clean'' alla ``raw'' (es. con Apache Web le
regole si scrivono nel file \texttt{.htaccess}), eventualmente diverse da un sito all'altro.
\begin{importante}[Rewrite engine e routing client-side]
Con un \emph{router lato client} (es. in React) il framework sostituisce la navigazione del
browser: al click su un link, se la URL \`e una ``location'' interna, configura la UI senza
vera navigazione (che resetterebbe l'app). Ma se l'utente \emph{scrive la URL nella barra},
il browser fa partire la richiesta al server: l'azione corretta \`e redirigerla sulla pagina
dell'applicazione -- e per questo serve il \textbf{Rewrite Engine}. In fase di sviluppo il
problema non si pone: il server HTTP che serve l'app client \`e interno all'ambiente di
sviluppo, dedicato solo ad essa, e sa cosa fare.
\end{importante}

\subsection{Logging}
Registrare informazioni sugli scambi HTTP, per scopi \emph{statistici} e di \emph{security}.

\section{Server HTTP e applicazione server-side: chi contiene chi?}

Un'applicazione server-side \`e un processo che resta in esecuzione, in attesa di processare
richieste HTTP ed elaborare risposte. Qual \`e la relazione con il server HTTP? Due casi
emblematici (architetture pi\`u complesse sono riconducibili a questi):

\begin{importante}[(a) Il Server HTTP ``contiene'' una o pi\`u app]
Il server esiste per s\'e, \`e installato e spesso fa altro oltre a dare accesso alle app
(spesso \`e un Web Server, e fra le risorse fornite c'\`e l'app client-side). Il server
contiene un \emph{motore di esecuzione} (es. una JVM) che tiene ``accese'' le applicazioni:
tipicamente crea un ambiente privato per ogni app (il \textbf{Web Container}) con una propria
istanza del motore, e un \emph{thread separato per ciascuna app} (modello
\textbf{multi-threaded}). Ogni app dichiara i pathname/route a cui risponde; il server li
registra in una \textbf{tabella di routing} e, all'arrivo di una richiesta sulla porta di
default, la inoltra all'app destinataria invocandone le funzioni, che restituiscono la
risposta da inviare al client.
\end{importante}

\begin{importante}[(b) L'app ``contiene'' il server HTTP]
Il punto d'ingresso \`e il motore di esecuzione del linguaggio, dentro il quale
l'applicazione \emph{si mette in ascolto} su una porta: le librerie del linguaggio offrono un
server HTTP di base che l'app istanzia e attiva. Il server interno ``avverte'' (meccanismo
event-driven) all'arrivo di una richiesta e la inoltra a una funzione dell'app, che
restituisce la risposta. La differenza chiave: in (a) \emph{il server ha il controllo e
lancia l'app}; in (b) \emph{l'app ha il controllo e lancia il server}. In questo modello
l'applicazione \`e un \textbf{Application Server basato su HTTP} (pu\`o anche fare da Web
Server, ma le relative funzionalit\`a vanno incluse esplicitamente, spesso via librerie). Si
abbina pi\`u facilmente a un server \emph{single-thread}; nulla vieta pi\`u app su porte
diverse dietro un \textbf{reverse proxy} sulla porta di default (avvicinando (b) ad (a)).
\end{importante}

\section{Caratteristiche generali di un'app server-side}

\textbf{Il mondo esterno.} Qualunque sia il modello, l'app comunica con l'esterno attraverso
il suo ambiente di esecuzione (che offre le librerie per farlo): il \emph{sistema operativo}
(in particolare il file system, spesso per logging); nel caso (a), \emph{informazioni di
contesto} fornite dal server; \emph{servizi o risorse di rete} -- fra cui la connessione a un
DBMS.

\textbf{Elementi e fasi di esecuzione:}
\begin{itemize}
  \item \textbf{Middleware}: prima che la richiesta arrivi all'app vera e propria, una o
        pi\`u funzioni effettuano \emph{elaborazioni preliminari}: la tipica funzione di
        middleware prende la richiesta, opzionalmente la trasforma e la ``passa oltre''
        (alla prossima funzione o all'app) \emph{oppure la blocca} -- per questo sono dette
        anche \textbf{filtri}. Usi: respingere richieste non autorizzate, interpretare e
        trasformare body/parametri, aggiungere informazioni. Utile per operazioni da eseguire
        \emph{sempre}, per ogni richiesta.
  \item \textbf{Individuazione della richiesta}: analizzare \emph{method + route} per
        determinare la funzione invocata, ed estrarre l'input (segmento di route parametrica,
        query string o body).
  \item \textbf{Elaborazione della risposta}: il ``core'' -- l'app esegue la funzionalit\`a e
        produce una risposta (risultati, o anche solo un acknowledgement).
  \item \textbf{Output della risposta}: la risposta viene inviata (tramite il server HTTP)
        \emph{in streaming}: la comunicazione verso il client \`e gestita come buffer, per
        inviare anche file grandi senza bloccare (anche lato client i dati arrivano in
        streaming: l'arrivo della risposta segna solo l'\emph{inizio} del flusso).
  \item \textbf{Errori}: richieste con operazioni non autorizzate, inconsistenze o parametri
        inaccettabili; oppure eccezioni nell'elaborazione. Va inviata una risposta con status
        code appropriato (\textbf{4xx} o \textbf{5xx} a seconda dei casi).
\end{itemize}

\section{Il ruolo dei framework server-side}

Il ruolo principale dei framework nelle app server-side \`e \textbf{automatizzare parte dei
compiti} e rendere pi\`u uniformi le interfacce verso il ``mondo esterno'': il framework si
frappone fra il core dell'applicazione e l'ambiente di esecuzione, facendosi carico di parte
dei compiti (quali, come, e come inserire codice personalizzato dipende dal framework). La
scelta del server HTTP \`e fortemente legata al \emph{linguaggio}, che determina
sostanzialmente il framework.

\textbf{Combo linguaggio/framework pi\`u usate nel web (2023):} PHP/Laravel, Python/Django,
Java/Spring, Python/Flask, Node.JS/Express.

\begin{esempio}[La strana storia di PHP]
Per decenni PHP \`e stato il principale linguaggio di back-end -- ma per un modello
architetturale diverso (pre-AJAX/JSON): il back-end \emph{generava dinamicamente pagine HTML}
e PHP semplifica molto questo compito. Esempi tipici: Wordpress e Moodle, i \textbf{CMS}
(content management system). Queste applicazioni hanno diffusione enorme e vengono ancora
aggiornate: PHP resta \emph{di gran lunga il linguaggio pi\`u usato nel Web} -- il che non
significa che sia il pi\`u richiesto o quello con cui si scriverebbe oggi un'app ex novo con
front-end molto interattivo (AJAX/JSON). Per applicativi aziendali su larga scala la
combinazione pi\`u richiesta oggi \`e \textbf{Java + Spring / Spring Boot}. Ha senso studiare
PHP oggi? S\`i, anche se non come primo linguaggio di web development.
\end{esempio}
